\documentclass[a4paper]{article}     
\usepackage{amsfonts}
\usepackage{amsmath}
\usepackage{amssymb}
\usepackage{graphicx}

\begin{document}
	
	\noindent \large Auteur: Victor Pena \\
	\noindent \large Studierichting: Informatica\\
	\noindent \large Oefeningenreeks 3F nr. 1.10\\
	
	\medskip
	
	\normalsize 
	
	$f(x)=x+\dfrac{1}{x} $\\
	
	$x+\dfrac{1}{x}= \dfrac{x^2+1}{x} \Rightarrow$ VA = $(x=0)$ \\
	
	$\lim\limits_{x\to+\infty} \dfrac{x^2+1}{x}	 = +\infty,\lim\limits_{x\to-\infty} \dfrac{x^2+1}{x}=-\infty \Rightarrow$ geen HA \\ 
	
	$\deg \left(x^2+1\right)=\deg \left(x\right)+1 \Rightarrow$ SA bestaat\\
	
	$\Rightarrow m= \lim\limits_{x\to\infty}\dfrac{f(x)}{x} = \lim\limits_{x\to\infty} \dfrac{x^2+1}{x^2} = 1 $\\
	
	$b = \lim\limits_{x\to\infty}f(x) -mx = \lim\limits_{x\to\infty} \dfrac{1}{x} = 0 \Rightarrow$ SA = $(y=x)$
	
	$d(x)= f(x)-x = \dfrac{1}{x}, x\to +\infty: d(x)>0,x\to -\infty: d(x)<0, d(x)$ is nooit gelijk aan $0$\\
	
	$\Rightarrow f(x)$ ligt onder het SA als $x\to-\infty$ en boven als $x\to+\infty$. $f(x)$ gaat naar $+\infty$ als $x\to0^+$ en gaat naar $-\infty$ als $x\to0^-$\\
	
	$\begin{tabular}{|c|ccccc|} \hline	
		$$&$-\infty$&$$&$0$&$$&$+\infty$ \\ \hline
		$d(x)$&$0$&$-$&$|$&$+$&$0$ \\
		$d'(x)$&$0$&$-$&$|$&$-$&$0$ \\
		$d''(x)$&$0$&$-$&$|$&$+$&$0$ \\ \hline
		$d(x)$&$0$&$\searrow$&$|$&$\searrow$&$0$ \\
		$d(x)$&$0$&$\frown$&$|$&$\smile$&$0$ \\ \hline
	\end{tabular}$\\
	
	
	
	\begin{thebibliography}{999}
		%\bibitem{a} Referentie
	\end{thebibliography}
\end{document} 
